\documentclass[11pt,a4paper]{ctexart}

% Fonts and layout
\usepackage{tgpagella}
\usepackage{geometry}
\geometry{
    a4paper,
    left=0.75in,
    right=0.75in,
    top=0.65in,
    bottom=0.65in
}

% Tables / lists / formatting
\usepackage{array}
\usepackage{tabularx}
\usepackage{enumitem}
\usepackage{titlesec}

% Links
\usepackage{hyperref}
\hypersetup{
    colorlinks=true,
    linkcolor=black,
    urlcolor=black,
    pdfborder={0 0 0}
}

% Section style
\titleformat{\section}{\fontfamily{phv}\large\bfseries}{}{0em}{}[\titlerule]
\titlespacing*{\section}{0pt}{0.65\baselineskip}{0.45\baselineskip}

% Paragraph/table spacing
\setlength{\parindent}{0pt}
\setlength{\parskip}{0.15em}
\setlength{\tabcolsep}{4pt}
\renewcommand{\arraystretch}{1.05}

% No page number
\pagestyle{empty}

% Column types
\newcolumntype{L}[1]{>{\raggedright\arraybackslash}p{#1}}
\newcolumntype{Y}{>{\raggedright\arraybackslash}X}

\begin{document}

% =========================
% Header
% =========================
\begin{center}
    {\fontfamily{phv}\selectfont\Large\bfseries Curriculum Vitae: 池卓倪 (Zhuoni Chi)}
\end{center}
\vspace{-0.35em}

\href{mailto:znchi@mit.edu}{znchi@mit.edu} \hfill \href{mailto:znchi@zju.edu.cn}{znchi@zju.edu.cn}\\
+1 (857) 331-9122 \hfill +86 18858185159

\vspace{0.2em}

% =========================
% Personal Information
% =========================
\section{Personal Information}
\begin{tabularx}{\textwidth}{L{1.4in}Y}
Birth   & 1999, Chongqing, China \\
Gender  & Male \\
Pronouns & he/him/his
\end{tabularx}

% =========================
% Education
% =========================
\section{Education}
\begin{tabularx}{\textwidth}{L{1.65in}Y}
2025--2026 &
\textbf{Visiting Student}, Massachusetts Institute of Technology \\
& Cambridge, MA, United States \\[0.15em]

2022--(2027 expected) &
\textbf{Ph.D. in Mathematics}, Zhejiang University \\
& Supervisor: Prof. Yifeng Liu \\
& Hangzhou, China \\[0.15em]

2018--2022 &
\textbf{B.S. in Mathematics}, Peking University \\
& Beijing, China \\
\end{tabularx}

% =========================
% Research Interests
% =========================
\section{Research Interests}
My research interests include algebraic number theory, automorphic representation theory, and their applications to the classical Langlands program. I have worked on arithmetic theta lifting for unitary groups, at the intersection of number theory, algebraic geometry, and representation theory.

I am also interested in formalizing mathematical olympiad problems and undergraduate-level mathematics in Lean4. I formalized about 50 AMC to USAMO level MO problems in Project Numina.

% =========================
% Research Work
% =========================
\section{Research Work}
\begin{enumerate}[label=(\arabic*), leftmargin=1.8em, itemsep=0.2em, topsep=0.15em]
    \item \textit{Spherical representations of unitary groups at ramified places and the arithmetic inner product formula} (preprint, available at \url{https://arxiv.org/abs/2602.02492})
\end{enumerate}

% =========================
% Academic Service
% =========================
\section{Academic Service}
\begin{itemize}[leftmargin=1.4em, itemsep=0.2em, topsep=0.15em]
    \item \textit{Teaching Assistant}: 2022 Homological Algebra, 2023 Algebra III at Zhejiang University; 2024 Lean4 Summer Camp at Peking University.
    \item \textit{Organizer} of five seminars on \'etale cohomology, arithmetic inner product formula, \(\mathrm{GL}_2\) local Langlands, \(p\)-adic Hodge theory, and \(p\)-adic geometry at Zhejiang University.
\end{itemize}

% =========================
% Honors & Awards
% =========================
\section{Honors \& Awards}
\begin{tabularx}{\textwidth}{L{1.1in}Y}
2017 &
\textbf{Gold Medal, 33rd Chinese Mathematical Olympiad (CMO), Rank 59/356} \\
& \textbf{Member of the National Training Team}, Hangzhou, China \\[0.15em]

2018 &
\textbf{Undergraduate Entry Scholarship}, Peking University, Beijing, China \\[0.15em]

2021 &
\textbf{Guoqin Scholarship}, Peking University, Beijing, China \\
\end{tabularx}

% =========================
% Skills
% =========================
\section{Skills}
\begin{tabularx}{\textwidth}{L{1.2in}Y}
Programming &
Python (no substantial software project experience yet);\\ &Lean4 (formal proof assistant; formalization of undergraduate mathematics and olympiad problems). \\[0.15em]


Languages &
Chinese (native), English (working proficiency). \\
\end{tabularx}

\end{document}
